\section{Labelled lifts: count the directions that
remain}

Let \(C\subset S^{D-1}\) be a kissing configuration. Split the larger
space as \(\mathbb R^D\oplus\mathbb R^k\). Choose pairwise distinct
nonzero tails \(t_i\) of norm less than one and subsets
\(A_i\subseteq C\). Write \(a_i=\sqrt{1-\|t_i\|^2}\) and
\(U=\bigcup_iA_i\). Retain \((x,0)\) for \(x\in C\setminus U\), insert
\((a_ix,t_i)\) for \(x\in A_i\), and add a pure-tail kissing code
\(Y\subset S^{k-1}\).

The count is

\[
|X|=|C|-|U|+\sum_i|A_i|+|Y|.
\]

Every original direction is deleted once, however many labels use it.
For distinct labelled pairs \((x,i)\ne(x',j)\) and pure tails
\(y\in Y\), the remaining conditions are

\[
a_ia_j\langle x,x'\rangle+\langle t_i,t_j\rangle\le\frac12,
\qquad
\langle t_i,y\rangle\le\frac12.
\]

These conditions separate tail geometry from head selection without
discarding their coupling.

\subsection{Multiplicity and the four-triangle
construction}

If every tail has squared norm \(h^2<1/2\), repeated copies of one
direction require

\[
\langle t_i,t_j\rangle\le h^2-\frac12.
\]

For \(m\) copies, the squared norm of the sum of their tails is at most

\[
\frac m2\bigl(1-m(1-2h^2)\bigr).
\]

Its nonnegativity implies \(m\le\lfloor(1-2h^2)^{-1}\rfloor\). At
\(h^2=1/3\), a direction can be used three times. Equality is realized
by a zero-sum triangle.

The twelve roots of \(D_3\), scaled to tail norm \(1/\sqrt3\), split
into four zero-sum triangles. Four disjoint 496-point Leech blocks, each
carried on one triangle, give

\[
196560-4(496)+12(496)+12=200540.
\]

The comparison allocation used multiplicities \((3,3,2,2,2)\). The new
allocation uses \((3,3,3,3)\), preserving 496 additional equatorial
points while keeping the same number of lifted points. The blocks
themselves are published PackingStar data; the change is their
assignment to the tail geometry.

For this fixed twelve-label system with at most \(\alpha\) directions
per label, the count is at most \(196560+8\alpha+12\). Thus the
construction attains the allocation capacity at \(\alpha=496\).

\subsection{The later 201,567-point
construction}

The later construction extends a public two-layer configuration. Its
first layer, taken unchanged from Lindow's 201,566-point contribution,
has 2,258 integer heads of squared norm 192. Of these, 2,090 delete one
Leech minimal vector each, and 168 delete none. Each head produces three
caps.

Use the integer Leech convention in which minimal vectors have squared
norm 32 and the kissing threshold is 16. A first-layer head \(y\)
produces

\[
(y/3,\;4v/\sqrt3),
\]

where \(v\) runs through its assigned \(D_3\) triangle. The exact head
thresholds are 48 on one triangle and 96 between different triangles.

A second-layer owner \(u\) is a Leech minimal vector assigned to one of
three coordinate lines. Replace its equatorial point by

\[
(\sqrt3\,u/2,\;\pm\sqrt8\,e_j).
\]

The supplied 311-owner witness has line sizes 105, 103, and 103. Every
first-layer head satisfies \(y\cdot u\le32\). Two owners on one line
have product at most eight; different lines have product at most
sixteen. Owners are distinct and do not overlap the first-layer deletion
set.

For example, the first/second-layer product is at most

\[
\frac{\sqrt3}{6}\,32+\frac{4\sqrt8}{\sqrt3}
=\frac{16\sqrt3+8\sqrt6}{3}<16.
\]

The second-layer/pure-tail product is at most \(\sqrt8\sqrt{32}=16\),
including equality. The two caps of one owner also have product sixteen.
These checks give

\[
196560-2090-311+3(2258)+2(311)+12=201567.
\]

Replacing the public input's 310 second-layer owners by the new
selection of 311 adds one point: each owner replaces one equatorial
point by two caps.
